\paragraph{幂基重整化协变性、Julia 动力学判别与阻尼阈值分相}\label{par:xi-discrete-abel-weil-powerbase-julia-threshold-laws}
本段在子小节 \ref{subsec:discrete-abel-weil-polar-rigidity} 与段落 \ref{def:abel-decimation-operators} 的记号下，
把“底数幂递归 \(b\mapsto b^m\)”统一为一条可审计链：
\[
\text{系数抽取协变}\ \Longrightarrow\ \text{极点幂映射}\ \Longrightarrow\ \text{Julia 分相判别}\ \Longrightarrow\ \text{Hardy 阈值二分}\ \Longrightarrow\ \sigma_* \text{ 的降采样反演}.
\]

\begin{theorem}[幂基重整化协变性与零点极点幂映射]\label{thm:xi-abel-powerbase-covariance-pole-power-map}
\leanverified{paper\_xi\_abel\_powerbase\_covariance\_pole\_power\_map}
固定 \(b>1\)，并沿用
\[
E_b(n)=b^{-n/2}\psi_0(b^n)-b^{n/2},\qquad
\mathcal A_b(r)=\sum_{n\ge 0}E_b(n)r^n
\]
与阻尼版
\[
\mathcal A_{b,\delta}(r):=\mathcal A_b(b^{-\delta}r)=\sum_{n\ge 0}E_b(n)b^{-\delta n}r^n\qquad(\delta\ge 0).
\]
对任意整数 \(m\ge 1\)，记 \(B:=b^m\)，并记抽取算子 \(\mathsf D_m\)（定义 \ref{def:abel-decimation-operators}）。
则有严格协变恒等式
\[
\mathcal A_{B,\delta}(r)=\mathsf D_m\!\bigl(\mathcal A_{b,\delta}\bigr)(r),
\qquad
H_{B,\delta}(r)=\mathsf D_m\!\bigl(H_{b,\delta}\bigr)(r),
\]
其中 \(H_{b,\delta}\) 为命题 \ref{prop:damped-poles} 的解析余项。
并且对任意非平凡零点 \(\rho\)，其阻尼极点
\[
r_{\rho,\delta}(b):=b^{-(\rho-(\frac12+\delta))}
\]
满足幂映射闭合律
\[
r_{\rho,\delta}(b^m)=\bigl(r_{\rho,\delta}(b)\bigr)^m.
\]
\end{theorem}

\begin{proof}
由定义直接计算
\[
E_{b^m}(n)
=(b^m)^{-n/2}\psi_0((b^m)^n)-(b^m)^{n/2}
=b^{-mn/2}\psi_0(b^{mn})-b^{mn/2}
=E_b(mn).
\]
因此
\[
\mathcal A_{b^m,\delta}(r)
=\sum_{n\ge 0}E_{b^m}(n)(b^m)^{-\delta n}r^n
=\sum_{n\ge 0}E_b(mn)b^{-m\delta n}r^n
=\mathsf D_m(\mathcal A_{b,\delta})(r).
\]
对余项
\(
H_{b,\delta}(r)=\sum_{n\ge 0}h_n r^n
\)
同理得到
\(
H_{b^m,\delta}(r)=\sum_{n\ge 0}h_{mn}r^n=\mathsf D_m(H_{b,\delta})(r)
\)。
最后由极点位置公式
\[
r_{\rho,\delta}(b^m)
=(b^m)^{-(\rho-(\frac12+\delta))}
=b^{-m(\rho-(\frac12+\delta))}
=\bigl(r_{\rho,\delta}(b)\bigr)^m
\]
即得。证毕。
\end{proof}

\begin{theorem}[极点动力学的 Julia 刚性与临界线判别]\label{thm:xi-abel-pole-julia-rigidity}
\leanverified{paper\_xi\_abel\_pole\_julia\_rigidity}
固定整数 \(m\ge 2\)，定义幂动力系统
\[
\Phi_m:\CC\setminus\{0\}\to\CC\setminus\{0\},\qquad \Phi_m(r)=r^m.
\]
对任意 \(\delta\ge 0\) 与任意非平凡零点 \(\rho\)，有等价链
\[
\Re\rho>\frac12+\delta
\Longleftrightarrow
|r_{\rho,\delta}(b)|<1
\Longleftrightarrow
\Phi_m^{\circ k}(r_{\rho,\delta}(b))\to 0,
\]
\[
\Re\rho<\frac12+\delta
\Longleftrightarrow
|r_{\rho,\delta}(b)|>1
\Longleftrightarrow
\Phi_m^{\circ k}(r_{\rho,\delta}(b))\to \infty,
\]
\[
\Re\rho=\frac12+\delta
\Longleftrightarrow
|r_{\rho,\delta}(b)|=1
\Longleftrightarrow
r_{\rho,\delta}(b)\in J(\Phi_m)=\{\,|r|=1\,\}.
\]
特别地，当 \(\delta=0\) 时
\[
\mathrm{RH}
\Longleftrightarrow
\forall\rho,\ \Re\rho=\frac12
\Longleftrightarrow
\{r_{\rho,0}(b)\}_{\rho}\subset J(\Phi_m).
\]
\end{theorem}

\begin{proof}
由命题 \ref{prop:damped-poles} 的模长公式
\[
|r_{\rho,\delta}(b)|
=b^{-(\Re\rho-(\frac12+\delta))}
\]
立即得到与 \(1\) 的大小比较。
另一方面，\(\Phi_m\) 迭代满足：\(|r|<1\Rightarrow r^{m^k}\to 0\)，\(|r|>1\Rightarrow r^{m^k}\to\infty\)，\(|r|=1\Rightarrow\) 轨道留在单位圆。
而 \(r\mapsto r^m\) 的 Julia 集是单位圆 \(\{|r|=1\}\)。
故全部等价成立；\(\delta=0\) 的特例即 RH 判别式。证毕。
\end{proof}

\begin{proposition}[重整化轨道的 Hardy 能量单调性与统一点态封口]\label{prop:xi-abel-hardy-energy-monotone-uniform-bound}
\leanverified{paper\_xi\_abel\_hardy\_energy\_monotone\_uniform\_bound}
取 \(\delta\ge 0\)，并假设 \(\mathcal A_{b,\delta}\in H^2(\mathbb D)\)。
则对任意 \(m\ge 1\) 有
\[
\|\mathcal A_{b^m,\delta}\|_{H^2(\mathbb D)}^2
=\sum_{n\ge 0}|E_b(mn)|^2 b^{-2m\delta n}
\le
\sum_{n\ge 0}|E_b(n)|^2 b^{-2\delta n}
=\|\mathcal A_{b,\delta}\|_{H^2(\mathbb D)}^2.
\]
并且对任意 \(0<r_0<1\) 成立统一点态界
\[
\sup_{m\ge 1}\ \sup_{|r|\le r_0}|\mathcal A_{b^m,\delta}(r)|
\le
\frac{\|\mathcal A_{b,\delta}\|_{H^2(\mathbb D)}}{\sqrt{1-r_0^2}}.
\]
\end{proposition}

\begin{proof}
由定理 \ref{thm:xi-abel-powerbase-covariance-pole-power-map}，
\(
\mathcal A_{b^m,\delta}=\mathsf D_m(\mathcal A_{b,\delta})
\)。
而 \(H^2\) 的系数刻画给出
\[
\|\mathsf D_m F\|_{H^2}^2=\sum_{n\ge 0}|a_{mn}|^2\le \sum_{n\ge 0}|a_n|^2=\|F\|_{H^2}^2
\]
（\(F(r)=\sum a_n r^n\)），得能量单调性。
点态估计由 Cauchy--Schwarz：
\[
|\mathsf D_mF(r)|
\le \Bigl(\sum_{n\ge 0}|a_{mn}|^2\Bigr)^{1/2}
\Bigl(\sum_{n\ge 0}|r|^{2n}\Bigr)^{1/2}
\le \frac{\|F\|_{H^2}}{\sqrt{1-|r|^2}}.
\]
取 \(F=\mathcal A_{b,\delta}\) 并在 \(|r|\le r_0\) 上取上确界即得。证毕。
\end{proof}

\begin{theorem}[阻尼阈值的稳定--失稳二分律]\label{thm:xi-abel-damping-threshold-bifurcation}
\leanverified{paper\_xi\_abel\_damping\_threshold\_bifurcation}
记
\[
\sigma_*:=\sup_{\rho}\Re\rho,\qquad
\delta_c:=\sigma_*-\frac12.
\]
固定 \(b>1\) 与 \(\delta\ge 0\)，考虑重整化轨道
\[
\{\mathcal A_{b^m,\delta}\}_{m\ge 1}.
\]
则有二分律：
\begin{enumerate}
  \item 若 \(\delta>\delta_c\)，则 \(\{\mathcal A_{b^m,\delta}\}_{m\ge 1}\) 在 \(\mathbb D\) 上局部一致有界，故为正规族；并且
  \[
  \sup_{m\ge 1}\|\mathcal A_{b^m,\delta}\|_{H^2(\mathbb D)}<\infty.
  \]
  \item 若 \(\delta<\delta_c\)，则对任意 \(r_0\in(0,1)\)，存在 \(m_0(r_0)\) 使得对所有 \(m\ge m_0(r_0)\)，
  \(\mathcal A_{b^m,\delta}\) 在 \(|r|\le r_0\) 内出现零点诱导极点；
  因而
  \[
  \sup_{|r|\le r_0}|\mathcal A_{b^m,\delta}(r)|=\infty.
  \]
\end{enumerate}
\end{theorem}

\begin{proof}
若 \(\delta>\delta_c\)，则存在 \(\varepsilon=\delta-\delta_c>0\) 使
\(
\Re\rho\le \frac12+\delta-\varepsilon
\)
对全部零点成立。
由定理 \ref{thm:hardy-criterion-delta} 得 \(\mathcal A_{b,\delta}\in H^2(\mathbb D)\)；
再由命题 \ref{prop:xi-abel-hardy-energy-monotone-uniform-bound} 得
\[
\sup_{m\ge 1}\sup_{|r|\le r_0}|\mathcal A_{b^m,\delta}(r)|
\le
\frac{\|\mathcal A_{b,\delta}\|_{H^2}}{\sqrt{1-r_0^2}}
\]
（任意 \(0<r_0<1\)），故局部一致有界并由 Montel 定理得到正规性。

若 \(\delta<\delta_c\)，则存在零点 \(\rho\) 使 \(\Re\rho>\frac12+\delta\)，即
\(
|r_{\rho,\delta}(b)|<1
\)。
由定理 \ref{thm:xi-abel-powerbase-covariance-pole-power-map}，
\[
|r_{\rho,\delta}(b^m)|
=|r_{\rho,\delta}(b)|^m
\xrightarrow[m\to\infty]{}0.
\]
故给定 \(r_0\in(0,1)\)，当 \(m\ge m_0(r_0)\) 时有
\(
|r_{\rho,\delta}(b^m)|<r_0
\)，即极点进入 \(|r|\le r_0\)。
于是 \(\sup_{|r|\le r_0}|\mathcal A_{b^m,\delta}(r)|=\infty\)。证毕。
\end{proof}

\begin{corollary}[任意固定步长降采样的增长率反演]\label{cor:xi-abel-fixed-step-growth-recovers-sigma-star}
\leanverified{paper\_xi\_abel\_fixed\_step\_growth\_recovers\_sigma\_star}
对任意整数 \(m\ge 1\)，定义
\[
\Lambda_{b,m}:=\limsup_{n\to\infty}\frac{1}{n}\log|E_b(mn)|.
\]
则恒有
\[
\Lambda_{b,m}=m\Bigl(\sigma_*-\frac12\Bigr)\log b,
\qquad
\sigma_*=\frac12+\frac{\Lambda_{b,m}}{m\log b}.
\]
即 \(\sigma_*\) 可由任意固定步长 \(m\) 的降采样序列指数增长率无损恢复。
\end{corollary}

\begin{proof}
由定理 \ref{thm:xi-abel-powerbase-covariance-pole-power-map}，\(E_{b^m}(n)=E_b(mn)\)。
再由定理 \ref{thm:radius-captures-sigma-star} 应用于底数 \(b^m\)：
\[
R_{b^m}=(b^m)^{-(\sigma_*-\frac12)}.
\]
由 Cauchy--Hadamard，
\[
\frac{1}{R_{b^m}}
=\limsup_{n\to\infty}|E_{b^m}(n)|^{1/n}
=\limsup_{n\to\infty}|E_b(mn)|^{1/n}.
\]
取对数即得
\[
\Lambda_{b,m}
=\log\frac{1}{R_{b^m}}
=\log\!\big((b^m)^{\sigma_*-\frac12}\big)
=m\Bigl(\sigma_*-\frac12\Bigr)\log b.
\]
整理得到反演式。证毕。
\end{proof}

\endinput
